\section{Радиальная огибающая и нормовый сертификат}
\label{bndnorm:section}\label{energy:section}
Положим $\mathcal A(v)=2H_0v+H_1v^2$. Для $P>q\ge0$ определим
\begin{equation}\label{boundarygap:time}
 \mathcal T(U;P,q)=\int_0^U
 \frac{dv}{\sqrt{\min\{q^2+\mathcal A(v),
                    P^2+\mathcal A(U)-\mathcal A(v)\}}}.
\end{equation}
При $P>0$, $0<r\le P/m$ зададим $A(P,r)$ уравнением
\begin{equation}\label{compact:radial-inverse}
 \mathcal T(A(P,r);P,P-mr)=r,\qquad A(P,0)=0.
\end{equation}

\begin{lemma}[Радиальное сравнение и точность]\label{compact:radial}
Уравнение \eqref{compact:radial-inverse} имеет единственное решение.
Пусть $u\in C^1([0,r])$, производная $u'$ абсолютно непрерывна и
\[
 u(0)=0,\quad u'(0)\ge0,\quad
 u'(s)\ge u'(0)+ms,\quad u'(r)\le P,\quad
 |u''(s)|\le H_0+H_1u(s)\quad\text{почти всюду}.
\]
Тогда $u(r)\le A(P,r)$. Равенство достигается на траектории с
начальным наклоном $P-mr$, конечным наклоном $P$ и кривизной
сначала $H_0+H_1u$, затем $-(H_0+H_1u)$.
Гладкие траектории с теми же крайними наклонами и константами
приближают это значение. При фиксированном $P$ функция $A(P,r)$
непрерывна и строго возрастает по $r$.

Если $H_1=0$, то
\begin{equation}\label{compact:radial-constant}
 A(P,r)=Pr+\beta r^2,\qquad
 \beta=\frac{H_0^2-2mH_0-m^2}{4H_0}.
\end{equation}
Если $H_1>0$, положим $k=\sqrt{H_1}$, $a=H_0/H_1$,
$T=kr$, $q=P-mr$ и
\[
 Y_b=a\cosh(T-b)+(q/k)\sinh(T-b),\qquad
 Q_b=ak\sinh(T-b)+q\cosh(T-b).
\]
На $[0,\min\{T,\pi/4\}]$ уравнение
\begin{equation}\label{compact:radial-switch}
 Q_b\cos b-kY_b\sin b=P
\end{equation}
имеет единственный внутренний корень, для которого
\begin{equation}\label{compact:radial-value}
 A(P,r)=Y_b\cos b+(Q_b/k)\sin b-a.
\end{equation}
\end{lemma}

\begin{proof}
Пусть $p=u'$, $U=u(r)$ и $q=P-mr$. Из условий следуют
$0\le p(0)\le q$ и $p(s)>0$ при $s>0$.
Дифференцирование $p^2-\mathcal A(u)$ и
$p^2+\mathcal A(u)$ почти всюду и интегрирование от двух концов дают
\begin{equation}\label{energy:two-sided}
 p(s)^2\le\min\{q^2+\mathcal A(u(s)),
                  P^2+\mathcal A(U)-\mathcal A(u(s))\}.
\end{equation}
Замена переменной $s$ на $u(s)$ даёт $\mathcal T(U;P,q)\le r$.
При $q=0$ особенность в нуле имеет порядок $v^{-1/2}$ и интегрируема.
Скачки второй производной этому рассуждению не мешают.

Проверим строгое обращение. Если $\mathcal A(U)\le P^2-q^2$,
в минимуме активно только первое выражение, и $\mathcal T$
строго возрастает. Иначе точка пересечения $v_c\in(0,U)$ задаётся
равенством $2\mathcal A(v_c)=P^2+\mathcal A(U)-q^2$.
При $H_1>0$ положим
\[
 z=U+a,\quad R^2=z^2+P^2/k^2,\quad
 w^2=\tfrac12(R^2+a^2-q^2/k^2),\quad w=v_c+a.
\]
Тогда
\[
 \mathcal T=
 \int_a^w\frac{dy}{\sqrt{k^2y^2+q^2-k^2a^2}}
 +\frac1k\{\arcsin(z/R)-\arcsin(w/R)\}.
\]
При дифференцировании слагаемые от точки пересечения сокращаются:
\begin{equation}\label{boundarygap:time-derivative}
 \mathcal T_U=\frac{P}{k^2R^2}
       +\frac{zw}{kR^2\sqrt{R^2-w^2}}>0.
\end{equation}
Ветви совпадают на границе. Первый интеграл неограниченно растёт
при $U\to\infty$, а второй неотрицателен.
При $H_1=0$
\[
 \mathcal T(U;P,q)=
 \begin{cases}
 (\sqrt{q^2+2H_0U}-q)/H_0,&2H_0U\le P^2-q^2,\\[1mm]
 (\sqrt{2q^2+2P^2+4H_0U}-q-P)/H_0,&2H_0U>P^2-q^2.
 \end{cases}
\]
Следовательно, $\mathcal T$ непрерывно и строго возрастает от нуля
до бесконечности. Это доказывает существование обратной функции
и верхнюю оценку.

При $H_1>0$ решим $u''=H_0+H_1u$ с $u(0)=0$, $u'(0)=q$
до времени $(T-b)/k$, затем $u''=-(H_0+H_1u)$ до $r$.
Конечный наклон $F(b)=Q_b\cos b-kY_b\sin b$ удовлетворяет
$F'(b)=-2kY_b\cos b<0$. При $b=0$
\[
 F(0)=ak\sinh T+q\cosh T>q+H_0r\ge P.
\]
На другом конце $F$ равно либо $q\cos T-ka\sin T<P$,
либо $(q-ka)e^{-(T-\pi/4)}/\sqrt2<P$.
Отсюда единственность внутреннего корня.
На первой дуге $u'\ge q+ms$, на второй производная убывает
до $P=q+mr$ и сохраняет ту же нижнюю границу.
Потеря неотрицательна, а обе энергетические оценки достигаются
на соответствующих дугах. Это даёт \eqref{compact:radial-value}.
При $H_1=0$ переключение от $H_0$ к $-H_0$ в точке
$r(H_0+m)/(2H_0)$ даёт \eqref{compact:radial-constant}.
При $H_0=m$ получается гладкая функция $qs+ms^2/2$.

В остальных случаях переключение внутреннее, и наклон в нём
строго больше $P$. Заменим знак кривизны гладким невозрастающим
переходом $\chi_{\epsilon,c}\in[-1,1]$ ширины $2\epsilon$ и решим
$u''=\chi_{\epsilon,c}(H_0+H_1u)$ с прежними начальными данными.
Непрерывная зависимость решения линейного ОДУ даёт равномерную
сходимость $u,u'$ при $\epsilon\downarrow0$.
Производная предельного конечного наклона по центру переключения
равна $2(H_0+H_1u(c))\cos(k(r-c))>0$ при $H_1>0$
и $2H_0$ при $H_1=0$. Поэтому небольшое изменение центра
сохраняет $u'(r)=P$ точно.
До перехода $u'\ge q+ms$; в переходе наклон остаётся выше $P$
с положительным запасом; после него убывает до $P$.
Все ограничения и крайние наклоны сохраняются.
Решения гладки в окрестности отрезка, совпадают с первой дугой
около нуля, а их потери стремятся к $A(P,r)$.

Пусть теперь $0<r_1<r_2\le P/m$, $d=r_2-r_1$, $q_i=P-mr_i$.
Перед почти экстремальной траекторией для $r_1$ добавим участок
$q_2t+mt^2/2$ на $[0,d]$. Его конечный наклон равен $q_1$,
а значение $a_0=q_2d+md^2/2>0$.
Продолжение $a_0+u(t-d)$ сохраняет нижнюю границу производной
и оценку кривизны. Производная склейки абсолютно непрерывна,
поэтому доказанное сравнение даёт
\begin{equation}\label{compact:radial-growth}
 A(P,r_2)\ge A(P,r_1)+q_2d+md^2/2>A(P,r_1).
\end{equation}
Непрерывность при $r>0$ следует из строгого уравнения
переключения, а при $r=0$ из верхней оценки с положительной
кривизной на всём пути. Вспомогательная склейка не меняет
допустимую область исходной задачи.
\end{proof}

Определим $\mathcal E(P)=A(P,P/m)$ при $P>0$, $\mathcal E(0)=0$.
При $H_1>0$ для неё достаточно одного скалярного корня; при $H_1=0$
\begin{equation}\label{energy:constant}
 \mathcal E(P)=\frac{P^2(H_0^2+2mH_0-m^2)}{4m^2H_0}.
\end{equation}

\begin{theorem}[Нормовый сертификат при сдвинутом RSI]
\label{bndnorm:main}\label{energy:main}
При условиях \eqref{boundarygap:assumptions} для каждого $x\in Q$
\begin{equation}\label{bndnorm:observable}
 \delta(x)\le\mathcal E(P(x)).
\end{equation}
При $x\ne x^*$ и радиальных наклонах $q=u'(0)$, $p=u'(r)$
\begin{equation}\label{bndnorm:radial-deficit}
 \delta(x)+\mathcal E(q)\le\mathcal E(p)\le\mathcal E(P(x)).
\end{equation}
Функция $\mathcal E$ непрерывна и для $0\le P_1<P_2$ удовлетворяет
\begin{equation}\label{compact:norm-growth}
 \mathcal E(P_2)\ge\mathcal E(P_1)+(P_2^2-P_1^2)/(2m).
\end{equation}
При каждом $P>0$ граница точна для наблюдения $P(x)=P$, в том
числе среди гладких примеров на отрезках с $f'(x^*)>0$.
\end{theorem}
\begin{proof}
Радиальное сведение и лемма~\ref{compact:radial} дают
$\delta(x)\le A(P(x),r)\le A(P(x),P(x)/m)$.
Продолжим стационарную экстремальную траекторию от $P_1/m$
до $P_2/m$ производной $ms$. Производная склейки абсолютно
непрерывна, а кривизна продолжения равна $m\le H_0+H_1u$.
То же сравнение доказывает \eqref{compact:norm-growth}.
При $P_1=0$ используем квадратичную траекторию.
Непрерывность следует из уравнения переключения; в нуле
применяется верхняя оценка с положительной кривизной.

Для радиальной поправки добавим перед исходным лучом гладкие
стационарные траектории на $[0,q/m]$ с конечным наклоном $q$
и потерями $a_j\to\mathcal E(q)$, затем продолжим как
$a_j+u(t-q/m)$. Первые производные совпадают, сдвинутое RSI
превращается в $v'(t)\ge mt$, а добавление $a_j\ge0$ сохраняет
оценку кривизны. Конечный наклон равен $p$, длина не больше $p/m$.
Сравнение даёт $a_j+\delta(x)\le\mathcal E(p)$.
Перейдём к пределу; при $q=0$ начальный участок пуст.

Точность при стационарном минимуме даёт лемма при $r=P/m$.
Чтобы получить положительный начальный наклон, возьмём
$q_j\downarrow0$ и стационарные почти экстремальные $v_j$
с конечным наклоном $P-q_j$ на $[0,(P-q_j)/m]$.
Функции $f_j(s)=v_j(s)+q_js$ сохраняют все константы,
имеют $f_j'(0)=q_j>0$, $f_j'(r_j)=P$ и
$f_j(r_j)\to\mathcal E(P)$.
\end{proof}

\begin{proposition}[Фиксированный начальный наклон]
\label{joint:fixed-slope}
При $P>q\ge0$ точный супремум радиальной потери с начальным
наклоном $q$, конечным наклоном $P$ и прежними ограничениями равен
\begin{equation}\label{joint:fixed-slope-time}
 \mathcal T(\mathcal E_q(P);P,q)=(P-q)/m.
\end{equation}
Его приближают гладкие траектории с теми же наклонами и константами.
При этом $\mathcal E_q(P)\le\mathcal E(P)-\mathcal E(q)$.
Если $H_1=0$, то
\begin{equation}\label{joint:fixed-slope-constant}
 \mathcal E_q(P)=\frac{P^2-q^2}{2m}
       +\frac{H_0^2-m^2}{4H_0m^2}(P-q)^2,
\end{equation}
а разность между прежней верхней границей и точной равна
\[
 \mathcal E(P)-\mathcal E(q)-\mathcal E_q(P)
 =\frac{q(P-q)(H_0^2-m^2)}{2H_0m^2}.
\]
\end{proposition}
\begin{proof}
Сдвинутое RSI даёт $r\le(P-q)/m$.
Энергетическое сравнение с точными крайними наклонами и строгое
возрастание $\mathcal T$ дают \eqref{joint:fixed-slope-time}.
Радиальная лемма при $r=(P-q)/m$ доказывает точность;
\eqref{bndnorm:radial-deficit} даёт оценку через разность.
При $H_1=0$ обращение явной формулы времени и вычитание
из \eqref{energy:constant} дают остальные равенства.
\end{proof}

\begin{corollary}[Полный зазор при стационарном минимуме]
\label{energy:gap}
Если $\nabla f(x^*)=0$, то для $W(g)=\mathcal E(\sqrt{mg})$
\begin{equation}\label{energy:gap-certificate}
 \delta(x)\le W\!\left(\min\{G(x),P(x)^2/m\}\right).
\end{equation}
Функция $W$ непрерывна, строго возрастает, не ограничена сверху
и удовлетворяет $W(g)\ge g/2$.
Граница точна при фиксированном полном зазоре.
\end{corollary}
\begin{proof}
При фиксированном $g>0$ определим $U(r)$ уравнением
$\mathcal T(U(r);g/r,0)=r$. Пусть $P=g/r$ и $v_c$ есть
точка пересечения энергетических границ. В смешанной ветви
\[
 \mathcal T+P\mathcal T_P
 =\int_0^{v_c}\frac{dv}{\sqrt{\mathcal A(v)}}
 +\int_{v_c}^{U}\frac{\mathcal A(U)-\mathcal A(v)}
 {[P^2+\mathcal A(U)-\mathcal A(v)]^{3/2}}\,dv>0.
\]
В ветви с одной границей $\mathcal T_P=0$.
Значит, $U'(r)=(\mathcal T+P\mathcal T_P)/(r\mathcal T_U)>0$.
Из $r\le\sqrt{g/m}$ получаем
$\delta(x)\le U(\sqrt{g/m})=W(g)$; объединим это с нормовой границей.
Свойства $W$ следуют из \eqref{compact:norm-growth}.
Стационарные примеры на отрезках имеют $G=P^2/m$ точно и
доказывают точность. Для этого вывода стационарность существенна.
\end{proof}
